\begin{abstract}
现有符号音乐生成研究较少直接处理后调性材料的显式控制，而使用受版权保护的当代作品进行可公开、可复现的评估也存在困难。本文提出一种乐谱层面的神经符号作曲辅助系统，其条件包括音级集合及音程向量、可选十二音序列及其变换、节奏轮廓、音乐姿态、乐器、声部数和乐谱跨度。合法的程序化生成器构建了 20,000 个训练片段、2,000 个验证片段和 2,000 个测试片段。系统采用六层因果 Transformer，在保留完整条件的事件窗口上训练，并在导出 MusicXML 和 JSON 报告之前，利用符号分析对采样候选进行排序。在固定的 2,000 条件测试集上，将候选数从 1 增至 4 并实施引导式选择后，音级集合覆盖率由 0.9963 升至 0.9999，循环序列顺序准确率由 0.1581 升至 0.2411，音程向量距离由 0.2390 降至 0.0045，节奏轮廓距离由 0.5113 降至 0.5041，姿态一致性由 0.4772 升至 0.4937。分别训练的单候选模型在移除各条件字段后，其对应自动指标均出现退化。13 个配置共尝试导出 260 份 MusicXML，全部通过结构解析、小节数和声部数检查，并整理了 20 个带分析报告的拟议模型样例。然而，所有神经配置的严格序列变换准确率均为 0，拟议模型的平均内容跨度比仅为 0.5195。因此，本文结论限定于合成测试床上的自动约束遵循，不涉及艺术质量或风格真实性。
\end{abstract}
